% Part: incompleteness
% Chapter: incompleteness-provability
% Section: first-incompleteness-thm

\documentclass[../../../include/open-logic-section]{subfiles}

\begin{document}

\olfileid{inc}{inp}{1in}

\olsection{पहिले अपूर्णता प्रमेय}

आता गोडेलने पहिल्या अपूर्णता प्रमेयासाठी दिलेल्या मूळ
सिद्धतेचे वर्णन करता येईल. अंकगणिताच्या भाषेचा विस्तार
असलेल्या भाषेतील $\Th{T}$ ही कोणतीही संगणनक्षमपणे
स्वयंसिद्धकीकृत उपपत्ती असू द्या, आणि $\Th{T}$ मध्ये
$\Th{Q}$ ची स्वयंसिद्धके समाविष्ट असू द्या. यावरून
विशेषतः $\Th{T}$ संगणनक्षम फलने आणि संबंध
निरूपित करते.

सूत्रे आणि सिद्धता संख्या म्हणून योग्य रीतीने सांकेतिक
केल्यास $\Prf[\Th{T}](x,y)$ हा संबंध संगणनक्षम
असतो, असा आपण युक्तिवाद केला आहे. येथे
$\Prf[\Th{T}](x,y)$ तेव्हाच खरा असतो जेव्हा
$x$ ही~$\Th{T}$ मध्ये गोडेल संख्या~$y$
असलेल्या !!{formula} च्या !!a{derivation}
ची गोडेल संख्या असते. गोडेलच्या मनात असलेल्या
विशिष्ट उपपत्तीसाठी हा संबंध आदिम पुनरावर्ती आहे
हेही त्याने आपल्या शोधपत्रातील 45 फलने आणि
संबंधांची यादी वापरून दाखवले. त्या यादीतील
45 वा संबंध, $x B y$, हा त्याने निवडलेल्या
$\Th{T}$ साठीचा $\Prf[\Th{T}](x,y)$
संबंधच आहे. गोडेल आपल्या शोधपत्रात जिथे
“पुनरावर्ती” हा शब्द वापरतो, तिथे आज आपण
“आदिम पुनरावर्ती” ही संज्ञा वापरू, हे लक्षात ठेवा.

$\Prf[\Th{T}](x,y)$ संगणनक्षम असल्यामुळे
तो $\Th{T}$ मध्ये निरूपणीय आहे. त्याचे
निरूपण करणाऱ्या सूत्रासाठी
$\OPrf[\Th{T}](x,y)$ हे चिन्ह वापरू.
$\OProv[\Th{T}](y)$ हे सूत्र
$\lexists[x][\OPrf[\Th{T}](x,y)]$
असे मानू. हे गोडेलच्या यादीतील 46 व्या
संबंधाचे, $ \fn{Bew}(y)$ चे, वर्णन करते.
गोडेलच्या नोंदीनुसार, हाच एक संबंध असा
आहे की तो “पुनरावर्ती आहे असे म्हणता येत
नाही.” त्याचा संभाव्य अर्थ असा: व्याख्येवरून
तो संगणनक्षम आहे हे स्पष्ट होत नाही; आणि
नंतरच्या प्रगतीतून तो प्रत्यक्षात संगणनक्षम
नाही हेही सिद्ध झाले.

$\Th{T}$ ही~$\Th{Q}$ समाविष्ट करणारी
!!{axiomatizable} उपपत्ती असू द्या. मग
$\Prf[\Th{T}](x, y)$ निर्णेय असतो आणि
म्हणून $\Th{Q}$ मध्ये
!!a{formula}~$\OPrf[\Th{T}](x, y)$
याने निरूपणीय असतो. वर वर्णन केल्याप्रमाणे
$\OProv[\Th{T}](y)$ हे सूत्र मानू.
स्थिरबिंदू उपपत्तीवरून असे सूत्र
$!G_\Th{T}$ मिळते की $\Th{Q}$
(आणि म्हणून $\Th{T}$ सुद्धा) पुढील
!!{derive}s करते:
\begin{equation}
\ollabel{eqn:qpf}
!G_\Th{T} \liff \lnot \OProv[\Th{T}](\gn{!G_\Th{T}}).
\end{equation}
लक्षात घ्या की $!G_\Th{T}$ चे
आशयतः विधान असे आहे: “$!G_\Th{T}$
हे~$\Th{T}$ मध्ये !!{derivable} नाही.”

\begin{lem}\ollabel{lem:cons-G-unprov}
$\Th{T}$ ही~$\Th{Q}$ चा विस्तार असलेली
सुसंगत आणि !!{axiomatizable} उपपत्ती
असेल, तर $\Th{T} \Proves/ !G_\Th{T}$.
\end{lem}

\begin{proof}
$\Th{T}$ हे $!G_\Th{T}$ ला
!!{derive}s करते असे मानू. मग
!!a{derivation} \emph{अस्तित्वात आहे}.
म्हणून एखाद्या संख्या $m$ साठी
$\Prf[\Th{T}](m,\Gn{!G_\Th{T}})$
हा संबंध खरा असतो. पण मग
$\Th{Q}$ हे
$\OPrf[\Th{T}](\num m, \gn{!G_\Th{T}})$
वाक्य !!{derive}s करते.
त्यामुळे $\Th{Q}$ हे
$\lexists[x][\OPrf[\Th{T}](x,\gn{!G_\Th{T}})]$
सुद्धा !!{derive}s करते; व्याख्येनुसार
हेच $\OProv[\Th{T}](\gn{!G_\Th{T}})$
आहे. \olref{eqn:qpf} वरून
$\Th{Q}$ हे $\lnot !G_\Th{T}$
!!{derive}s करते. $\Th{T}$ हा
$\Th{Q}$ चा विस्तार असल्यामुळे
$\Th{T}$ सुद्धा तसेच करते.
म्हणजे $\Th{T}$ हे $!G_\Th{T}$
!!{derive}s करत असेल, तर
$\lnot !G_\Th{T}$ सुद्धा
!!{derive}s करते, असे आपण दाखवले.
त्यामुळे ते विसंगत ठरेल.
\end{proof}

\begin{defn}
\ollabel{thm:oconsis-q}
उपपत्ती $\Th{T}$ ला
\emph{$\omega$-सुसंगत} म्हणतात, जर
पुढील अट पूर्ण होत असेल: कोणतेही
वाक्य $\lexists[x][!A(x)]$ घेऊ.
$\Th{T}$ हे $\lnot !A(\num 0)$,
$\lnot !A(\num 1)$,
$\lnot !A(\num 2)$, \dots
!!{derive}s करत असेल, तर
$\Th{T}$ हे $\lexists[x][!A(x)]$ सिद्ध करत नाही.
\end{defn}

प्रत्येक $\omega$-सुसंगत उपपत्ती
सुसंगतसुद्धा असते, हे लक्षात घ्या.
कारण $\Th{T}$ विसंगत असेल, तर
प्रत्येक~$!A$ साठी $\Th{T} \Proves !A$
असते. विशेषतः $\Th{T}$ विसंगत
असेल, तर प्रत्येक~$n$ साठी
$\lnot !A(\num n)$ हे ते
!!{derive}s करते आणि
$\lexists[x][!A(x)]$ हेही
!!{derive}s करते.
म्हणून विसंगत $\Th{T}$ हे
$\omega$-विसंगत असते. याच्या
प्रतिवर्ती विधानावरून
$\omega$-सुसंगत $\Th{T}$
सुसंगत असलेच पाहिजे.

\begin{lem}\ollabel{lem:omega-cons-G-unref}
$\Th{T}$ ही~$\Th{Q}$ चा विस्तार
असलेली $\omega$-सुसंगत,
!!{axiomatizable} उपपत्ती
असेल, तर
$\Th{T} \Proves/ \lnot !G_\Th{T}$.
\end{lem}

\begin{proof}
$\Th{T}$ हे $\lnot !G_\Th{T}$
!!{derive}s करत असेल, तर ते
$\omega$-विसंगत असते, हे दाखवू.
$\Th{T}$ हे $\lnot !G_\Th{T}$
!!{derive}s करते असे मानू.
$\Th{T}$ विसंगत असेल, तर ते
$\omega$-विसंगत आहेच.
नसेल, तर $\Th{T}$ सुसंगत आहे.
म्हणून \olref{lem:cons-G-unprov}
वरून ते $!G_\Th{T}$ ला
!!{derive} करत नाही. $\Th{T}$
मध्ये $!G_\Th{T}$ ची
!!{derivation} नसल्यामुळे
$\Th{Q}$ पुढील विधाने
!!{derive}s करते:
\[
\lnot \OPrf[\Th{T}](\num 0, \gn{!G_\Th{T}}), \lnot \OPrf[\Th{T}](\num 1,
\gn{!G_\Th{T}}), \lnot \OPrf[\Th{T}](\num 2, \gn{!G_\Th{T}}), \dots
\]
त्यामुळे $\Th{T}$ सुद्धा ती
सिद्ध करते. दुसरीकडे
\olref{eqn:qpf} वरून
$\lnot !G_\Th{T}$ हे
$\lexists[x][\OPrf[\Th{T}](x,\gn{!G_\Th{T}})]$
याच्याशी सममूल्य असते.
म्हणून $\Th{T}$ हे
$\omega$-विसंगत आहे.
\end{proof}

\begin{prob}
  प्रत्येक $\omega$-सुसंगत उपपत्ती
  सुसंगत असते. याचे उलट खरे नसते
  हे दाखवा: सुसंगत पण
  $\omega$-विसंगत उपपत्ती असतात.
  यासाठी $\Th{Q} \cup
  \{\lnot !G_\Th{Q}\}$ हा संच
  सुसंगत पण $\omega$-विसंगत
  आहे असे दाखवा.
\end{prob}

\begin{thm}
\ollabel{thm:first-incompleteness}
$\Th{T}$ ही~$\Th{Q}$ चा
$\omega$-सुसंगत,
!!{axiomatizable} विस्तार असलेली
कोणतीही उपपत्ती असू द्या.
मग $\Th{T}$ संपूर्ण नाही.
\end{thm}

\begin{proof}
  $\Th{T}$ हे $\omega$-सुसंगत
  असेल, तर ते सुसंगत असते.
  म्हणून \olref{lem:cons-G-unprov}
  वरून $\Th{T}
  \Proves/ !G_\Th{T}$.
  \olref{lem:omega-cons-G-unref}
  वरून $\Th{T} \Proves/
  \lnot !G_\Th{T}$. म्हणजे
  $\Th{T}$ अपूर्ण आहे: ते
  $!G_\Th{T}$ किंवा
  $\lnot !G_\Th{T}$ यांपैकी
  एकही विधान !!{derive}s करत नाही.
\end{proof}

\end{document}
